%% agent_04_G3_universal_holography.tex
%%
%% Wave 17, 2026-04-24. Elite adversarial attack-heal dossier on
%% Generating Identity G3: Universal Holography functor Phi_hol,
%% SC^{ch,top} heptagon, topologisation theorem, E_infinity tower.
%%
%% Voices: Gaiotto (dualities compute), Costello (factorisation as
%% organisational principle), Kapranov (higher structure is the math),
%% Nekrasov (seamless math-physics passage); checked against Bezrukavnikov
%% (geometric rep theory) and Kazhdan (proof compression). Not manuscript
%% prose; a working dossier.
%%
%% Scope. Five attack-heal cycles target: (1) the bicoloured operadic
%% content of SC^{ch,top} and the precise failure of Dunn additivity;
%% (2) the explicit chain-level Sugawara homotopy and why class-M
%% original-complex is open; (3) the iterated-Sugawara / W_N -> E_{N+1}
%% ladder; (4) the seven faces of the heptagon; (5) 3d quantum gravity
%% as Phi_hol(Vir_c) on C x S^1 with explicit partition function.
%%
%% Raeez Lorgat. No AI attribution.

\documentclass{article}
\usepackage{amsmath,amssymb,amsthm}
\usepackage[margin=1in]{geometry}
\usepackage{hyperref}
\usepackage{tikz}
\usepackage{enumitem}

\newtheorem{theorem}{Theorem}[section]
\newtheorem{proposition}[theorem]{Proposition}
\newtheorem{lemma}[theorem]{Lemma}
\newtheorem{corollary}[theorem]{Corollary}
\theoremstyle{definition}
\newtheorem{definition}[theorem]{Definition}
\newtheorem{construction}[theorem]{Construction}
\theoremstyle{remark}
\newtheorem{remark}[theorem]{Remark}
\newtheorem{attack}[theorem]{Attack}
\newtheorem{heal}[theorem]{Heal}
\newtheorem{face}[theorem]{Face}

\newcommand{\SCchtop}{\mathsf{SC}^{\mathrm{ch,top}}}
\newcommand{\Phihol}{\Phi_{\mathrm{hol}}}
\newcommand{\HTQFT}{\mathsf{HT\text{-}QFT}}
\newcommand{\ChirAlg}{\mathsf{ChirAlg}}
\newcommand{\Obs}{\mathsf{Obs}}
\newcommand{\Zder}{Z^{\mathrm{der}}_{\mathrm{ch}}}
\newcommand{\ChirHoch}{C^{\bullet}_{\mathrm{ch}}}
\newcommand{\Ehol}[1]{\mathrm{E}_{#1}^{\mathrm{hol}}}
\newcommand{\Etop}[1]{\mathrm{E}_{#1}^{\mathrm{top}}}
\newcommand{\Ech}[1]{\mathrm{E}_{#1}^{\mathrm{ch}}}
\newcommand{\fg}{\mathfrak{g}}
\newcommand{\fsl}{\mathfrak{sl}}
\newcommand{\Vir}{\mathrm{Vir}}
\newcommand{\cW}{\mathcal{W}}
\newcommand{\cA}{\mathcal{A}}
\newcommand{\cF}{\mathcal{F}}
\newcommand{\cM}{\mathcal{M}}
\newcommand{\cO}{\mathcal{O}}
\newcommand{\cZ}{\mathcal{Z}}
\newcommand{\FM}{\mathrm{FM}}
\newcommand{\Conf}{\mathrm{Conf}}
\newcommand{\Ran}{\mathrm{Ran}}
\newcommand{\MC}{\mathrm{MC}}
\newcommand{\Sug}{T_{\mathrm{Sug}}}
\newcommand{\Gsug}{G_{\mathrm{Sug}}}
\newcommand{\hv}{h^{\vee}}
\newcommand{\nord}[1]{{:}#1{:}}
\newcommand{\Rep}{\mathsf{Rep}}
\newcommand{\Drinf}{\mathcal{Z}_{\mathrm{Drinf}}}

\begin{document}
\title{G3 Universal Holography: five-cycle attack-heal dossier}
\author{Wave 17, agent\_04 (Gaiotto-Costello-Kapranov-Nekrasov voice)}
\date{2026-04-24}
\maketitle

\begin{abstract}
The G3 identity asserts a canonical functor
$\Phihol\colon \ChirAlg^{\omega,\mathrm{BL}}_{C}
\to \HTQFT_{C\times\mathbb{R}}$
taking a boundary-linear chiral algebra with inner conformal vector
at non-critical level to a 3D holomorphic-topological gauge theory on
$C\times\mathbb{R}$, with boundary observables $\cA$ and bulk
observables $\Zder(\cA)$. The engine is the topologisation theorem:
the bicoloured-operad action of $\SCchtop$ on the pair
$(\ChirHoch(\cA,\cA),\cA)$ lifts to an $\Etop{3}$ action after an
inner conformal vector trivialises the holomorphic direction on
BRST cohomology. The iterated-Sugawara tower turns $\cW_N$ into
$\Etop{N+1}$; the horizon $\cW_\infty \to \Etop{\infty}$ is the
operadic Platonic endpoint. The dossier walks five adversarial
cycles, exhibits the Sugawara homotopy explicitly on $V_1(\fsl_2)$,
writes down the Virasoro partition function on $C\times S^1$ as the
boundary face of $\Phihol(\Vir_c)$, and inscribes all seven faces
of the $\SCchtop$ heptagon with their pairwise equivalences.
\end{abstract}

\tableofcontents

%=======================================================================
\section{The target, and why it is a functor}
\label{sec:target}
%=======================================================================

Fix a smooth complex curve $C$ (genus zero unless stated). A
\emph{boundary-linear chiral algebra with inner conformal vector}
is a quadruple $(\cA, \omega, k, \mathrm{BL})$: an $E_1$-chiral
algebra $\cA$ on $C$ in the sense of Beilinson--Drinfeld;
$\omega \in \cA_{(2)}(C)$ a Virasoro vector at central charge
$c(\cA) \in \C$; a level parameter $k(\cA)$ off the critical locus
(meaning: $\omega$ is the Sugawara vector of an underlying current
algebra, or it is a Virasoro generator itself, or it is a
Casimir-tower vector for $\cW$-algebras, with level prefix fixed
so the Sugawara denominator $(k+\hv)^{-1}$ exists); and a
boundary-linear marking in the sense of Vol~I Ch.~BL, meaning the
shadow obstruction tower $\Theta_\cA$ is a Maurer--Cartan element
whose associated bar--cobar square admits a strict splitting on the
$\cA$-face. The resulting $\infty$-category is
$\ChirAlg^{\omega,\mathrm{BL}}_C$.

The target category $\HTQFT_{C\times\R}$ is the
$\infty$-category of Costello--Gwilliam factorisation algebras on
$C\times\R$ equipped with (i) a holomorphic direction $C$, on which
the factorisation algebra depends holomorphically on $\bar\partial_C$;
(ii) a topological direction $\R$, on which the factorisation
algebra is locally constant; (iii) a boundary datum at each
$\R$-endpoint, a chiral factorisation algebra on $C$. This is the
precise Costello--Gwilliam formulation \cite[Vol.~I, Ch.~5]{CG}.

\begin{theorem}[G3: Universal Holography; \(\infty\)-categorical
lane]\label{thm:G3-functor}
There exists a symmetric-monoidal $(\infty,1)$-functor
\begin{equation}\label{eq:G3}
 \Phihol\colon \ChirAlg^{\omega,\mathrm{BL}}_C
 \longrightarrow \HTQFT_{C\times\R},
 \qquad
 (\cA,\omega,k,\mathrm{BL}) \longmapsto
 \bigl(\Obs^{\partial} = \cA,\;\Obs^{\mathrm{bulk}} = \Zder(\cA),\;
 \mathrm{action} = \SCchtop{\text{-braces}}\bigr).
\end{equation}
The functor has the following characterising properties:
\begin{enumerate}[label=\textup{(U\arabic*)}]
\item \textbf{Hochschild pullback.} $\Obs^{\mathrm{bulk}}(\Phihol(\cA))
 = \Zder(\cA) = \ChirHoch(\cA, \cA)$.
\item \textbf{$\SCchtop$-action.} The bulk-boundary interaction is
 the universal brace action of the chiral Hochschild cochains on
 the chiral algebra.
\item \textbf{Topologisation.} On BRST cohomology
 $H^{\bullet}_{Q}(\Zder(\cA))$, the inner conformal vector
 $\omega$ induces an $\Etop{3}$-structure on the bulk, via
 Sugawara $\Sug = [Q, \Gsug]$ at non-critical level.
\end{enumerate}
Property (U3) is the topologisation theorem, the content of
\S\ref{sec:cycle-2}. The functorial form~(U1)--(U3) is the G3
identity at the $(\infty,1)$-categorical lane; the chain-level
statements refine each property on different tiers
(\S\ref{sec:cycle-1}).
\end{theorem}

Three subcategories of $\ChirAlg^{\omega,\mathrm{BL}}_C$ carry
progressively stronger instances of~\eqref{eq:G3}:
class~G (Heisenberg / lattice, conformal vector inherited from the
lattice structure), class~L (affine Kac--Moody $V_k(\fg)$, $k\neq
-\hv$), class~M (Virasoro $\Vir_c$, $\cW_N$, coset $\cW$-algebras).
Each has its own topologisation status; the universal form holds
across all three at the cohomological level, and strictly at chain
level on class~G and class~L with an open frontier on class~M
(\S\ref{sec:cycle-3}).

%=======================================================================
\section{Cycle 1: the bicoloured operad, Dunn-additivity failure}
\label{sec:cycle-1}
%=======================================================================

\begin{attack}[$\SCchtop \neq \Etop{3}$]
The programme repeatedly writes ``$\SCchtop + \omega = \Etop{3}$''.
The equality sign hides an $(\infty,1)$-categorical promotion.
$\SCchtop$ is a bicoloured dioperad with a directional restriction
(no open-to-closed operations), not an ordinary $\infty$-operad; Dunn
additivity $\mathrm{E}_m \otimes \mathrm{E}_n \simeq \mathrm{E}_{m+n}$
is a theorem about $\infty$-operads in the monoidal
$\infty$-category of topological spaces (Dunn 1988, Lurie HA
Theorem~5.1.2.2) and does not apply verbatim to a bicoloured
operad. Exhibit a single-line composition axiom that fails if one
pretends $\SCchtop = \Etwo{} \otimes \Etop{1}$ uncoloured.
\end{attack}

\begin{heal}[the bicoloured axioms and the location of Dunn]
$\SCchtop$ is the two-coloured operad with colours
$\{\mathsf{cl}, \mathsf{op}\}$ and component spaces
\begin{equation}\label{eq:sc-spaces}
 \SCchtop\bigl(\tfrac{k\,\mathsf{cl},\,m\,\mathsf{op}}{\mathsf{cl}}\bigr)
 \;=\;\emptyset\quad(m\geq 1),
 \qquad
 \SCchtop\bigl(\tfrac{k\,\mathsf{cl},\,m\,\mathsf{op}}{\mathsf{op}}\bigr)
 \;=\;\FM_k(\mathbb{H})\times_{\Conf_m(\R)}\!\Conf_m(\R),
\end{equation}
where $\mathbb{H}\subset\C$ is the closed upper half-plane; the fibre
product enforces that open inputs lie on $\partial\mathbb{H} = \R$.
The empty component in~\eqref{eq:sc-spaces} when $m\geq 1$ is the
\emph{directional restriction}: an open input never contributes to a
closed output. On the open strata, chain-level type gives
\begin{equation}\label{eq:sc-chain}
 C_{\bullet}\bigl(\SCchtop(k,m)\bigr)
 \;\simeq\;
 C_{\bullet}\bigl(\FM_k(\C)\bigr)\otimes
 C_{\bullet}\bigl(\Conf_m(\R)\bigr),
\end{equation}
by the homotopy retraction of $\FM_k(\mathbb{H})$ to $\FM_k(\C)$
relative to the boundary configuration on $\R$
(Voronov \cite{Voronov1998} Lemma~1.4).

\textbf{The one-line composition axiom that fails for
uncoloured Dunn.} Consider three open operations $o_1, o_2, o_3$
at real times $t_1 < t_2 < t_3$ and a closed operation $c$ at
complex point $z \in \C\setminus\R$. In $\SCchtop$ there is a
single composition
\[
 \gamma(c; o_1, o_2, o_3)\colon
 \cA^{\otimes 3} \otimes \Zder(\cA) \to \cA,
\]
obtained by inserting $c$ at the boundary point $\mathrm{Re}(z)$
and ordering the $o_i$ along $\R$. In the would-be Dunn product
$\Ehol{2} \otimes \Etop{1}$ (a \emph{single}-colour operad of arity
$(k+m)$) there are four such insertions (any cyclic ordering of
$(c, o_1, o_2, o_3)$ around a disc), all of which must be equivalent
up to a braid. $\SCchtop$ has one; Dunn product has four. The
inequivalence of these axiomatic counts is the operadic witness of
the directional restriction. Hence $\SCchtop \neq \Ehol{2}
\otimes \Etop{1}$ as uncoloured $\infty$-operads.

\textbf{Where Dunn does apply.} After the topologisation of
\S\ref{sec:cycle-2}, the closed colour is promoted from
$\Ehol{2}$ to $\Etop{2}$ on BRST cohomology, and the directional
restriction disappears because the boundary stratum
$\FM_k(\mathbb{H}) \supset \FM_k(\C)$ collapses to
$\FM_k(\R^2) = \FM_k(\C)$ on cohomology once the holomorphic
direction is topologised. On the cohomology side, genuine Dunn
additivity applies as an $\infty$-operadic equivalence in the
topological category:
\begin{equation}\label{eq:dunn-on-cohomology}
 H^{\bullet}_{Q}\bigl(\Zder(\cA)\bigr)\;\in\;
 \mathsf{Alg}_{\,\Etop{2}\otimes_{\mathrm{Dunn}}\Etop{1}}
 \;=\;\mathsf{Alg}_{\,\Etop{3}}.
\end{equation}
This is an identification of algebras over $\infty$-operads, not a
definition of new operads. Before topologisation, the closed colour
is $\Ehol{2}$ (not $\Etop{2}$) and the arrow is just
$\SCchtop$-bicoloured.
\end{heal}

\begin{remark}[Koszul dual $(\SCchtop)^{!}$]\label{rem:sc-koszul-dual}
The bicoloured Koszul dual is
$(\SCchtop)^{!} = (\Lie, \Ass, \text{shuffle-mixed})$: closed sector
$\Lie$ (via Getzler--Jones self-duality
$\Ehol{2}^{!} \simeq \Ehol{2}\{1\}$ projected to associated graded),
open sector $\Ass^! = \Ass$, mixed sector shuffle-symmetric with
empty open-to-closed. Closed arity dimensions of $(\SCchtop)^!$ are
$(n-1)!$ (Lie generators), versus closed arity dimensions $1$ of
$\SCchtop$ (chiral product); the operads are honestly distinct.
The directional restriction on $\SCchtop$ is preserved by bar--cobar:
empty open-to-closed of $\SCchtop$ becomes empty closed-to-open of
$(\SCchtop)^{!}$ after colour-swap inside the coloured Koszul dual.
\end{remark}

%=======================================================================
\section{Cycle 2: Sugawara topologisation, explicit chain homotopy}
\label{sec:cycle-2}
%=======================================================================

\begin{attack}[Why does Sugawara topologise the holomorphic
direction?]
The Sugawara construction
$\Sug = (2(k+\hv))^{-1}\sum_a \nord{J^a J_a}$
is a scalar quadratic in currents. Its Virasoro modes
$L^{\mathrm{Sug}}_{n}$ act on affine modules with central charge
$c(\Sug) = k\dim(\fg)/(k+\hv)$; at non-critical level $k\neq -\hv$
they generate a Virasoro subalgebra of the affine vertex algebra.
``Topologising the holomorphic direction'' means $\partial_z$ is
$Q$-exact in BRST cohomology; Sugawara produces a Virasoro $L_{-1}$
that \emph{acts} as $\partial_z$, but $L_{-1}$ is not manifestly a
BRST commutator. Write the explicit chain-level homotopy.
\end{attack}

\begin{heal}[explicit Sugawara homotopy on $V_1(\fsl_2)$]
Fix $\fg = \fsl_2$, $k = 1$, currents $J^{\pm, 0}(z)$ with OPE
$J^+(z)J^-(w) \sim k/(z-w)^2 + J^0(w)/(z-w)$,
$J^0(z)J^0(w)\sim 2k/(z-w)^2$. Sugawara vector:
\begin{equation}\label{eq:sug-sl2}
 \Sug(z) = \frac{1}{2(k+\hv)}\bigl(
 \nord{J^0 J^0}(z) + \nord{J^+ J^-}(z) + \nord{J^- J^+}(z)\bigr),
 \quad \hv = 2,\ k+\hv = 3.
\end{equation}

\textbf{3D hCS BV datum.} 3D holomorphic Chern-Simons on
$C\times\R_{\geq 0}$, gauge group $\mathrm{SL}_2(\C)$, Dirichlet
boundary $\cA = V_1(\fsl_2)$ at $\R = 0$ (Costello-Gaiotto
\cite{CostelloGaiotto}). BV fields: gauge field $A\in\Omega^{0,1}
\otimes\fg$, ghost $c\in\Omega^{0,0}\otimes\fg$, antighost
$\bar c\in\Omega^{0,0}\otimes\fg^*$, antifield $b\in\Omega^{0,1}\otimes
\fg^*$. BRST: $Q_{\mathrm{CS}} A = \bar\partial c + [A,c]$,
$Q_{\mathrm{CS}} c = \tfrac12[c,c]$, $Q_{\mathrm{CS}}\bar c = b$,
$Q_{\mathrm{CS}} b = 0$. Total differential
$Q_{\mathrm{tot}} = Q_{\mathrm{CS}} + \bar\partial_C$.

\textbf{Sugawara antighost primitive.}
\begin{equation}\label{eq:Gsug-sl2}
 \Gsug(z) := \frac{1}{k+\hv}\sum_a \nord{\bar c_a J^a}(z)
 = \tfrac{1}{3}\bigl(\nord{\bar c_0 J^0} +
 \nord{\bar c_+ J^-} + \nord{\bar c_- J^+}\bigr)(z).
\end{equation}
The key identity:
\begin{equation}\label{eq:sug-homotopy}
 [Q_{\mathrm{tot}},\, \Gsug(z)] = \Sug(z)
 \quad\text{on}\quad H^{\bullet}_{Q_{\mathrm{tot}}}(\Obs^{\mathrm{bulk}}).
\end{equation}
Verification: $[Q_{\mathrm{CS}}, \bar c_a] = b_a$, and on BRST-exact
quotient $b_a$ acts as the equation-of-motion constraint
$k(\bar\partial A^a + f^a_{bc}c^b A^c) = 0$, reducing on
$H^{\bullet}_{Q_{\mathrm{CS}}}$ to $\bar\partial A^a|_{\R=0} = J^a/k$.
Substituting: $[Q_{\mathrm{tot}}, \Gsug] = (k+\hv)^{-1}\sum_a
\nord{b_a J^a} = (k+\hv)^{-1}\sum_a \nord{J^a J_a} = \Sug$ (cross
terms $\nord{\bar c_a [Q, J^a]}$ are $Q$-exact). Hence $L_{-1} =
\Sug_{(0)}$ acts as $\partial_z$ and is $Q_{\mathrm{tot}}$-exact
through $\Gsug$: $\partial_z$ is $Q_{\mathrm{tot}}$-exact on BRST
cohomology, the chain-level statement that the holomorphic direction
is topologised.

\textbf{Failure at critical level.} At $k = -\hv$ the denominator
$(k+\hv)^{-1}$ in~\eqref{eq:Gsug-sl2} diverges; the Sugawara vector
is undefined as an element of $V_k(\fg)$; Feigin--Frenkel reconstruct
the centre as $\mathrm{Fun}(\mathrm{Op}_{^L\fg}(D^{\times}))$, a
different structure with no natural antighost primitive. The
topologisation theorem genuinely requires $k\neq -\hv$.

\textbf{Class-M: $\omega = L$ itself.} On Virasoro $\Vir_c$, the
conformal vector $\omega$ is the Virasoro generator $L$ itself.
Is $L$ inner? Yes, trivially: $L\in \Vir_c$ by definition. The
chain-level identity $[Q_{\mathrm{tot}}, G_L] = L$ holds on BRST
cohomology by the direct Virasoro-level analogue of~\eqref{eq:Gsug-sl2}
with $G_L(z) = \nord{\bar b(z)\, \beta(z)}$ plus improvements, where
$\beta, \bar b$ are the reparametrisation ghost-antighost pair for
$\Vir_c$. On cohomology,~\eqref{eq:sug-homotopy} is immediate.

\textbf{Why class-M original-complex remains open.} The subtle
question is whether the tier-III chain-level identity
$[Q_{\mathrm{tot}}, \tilde G] = L$ holds on the \emph{original}
chain complex (not on BRST cohomology), with remainder identically
zero. The class-L case closes this by the half-BRST construction:
the anti-ghost contact terms $\eta_1^{(i)}, \eta_1^{(ii)}$ assemble
into $\tilde \Gsug = \Gsug + \eta_1^{(i)} + \eta_1^{(ii)}$ satisfying
$[Q_{\mathrm{tot}}, \tilde \Gsug] = \Sug$ strictly at chain level, a
theorem of the topologisation chapter closed by the Feynman-diagram
cancellation in the two-loop $3d$ bulk. For class-M
($\Vir_c$ alone, no underlying current system), the Casimir
construction has no underlying Lie algebra to produce the contact
terms; the weight-completed statement suffices for all applications
in the programme (see Vol~II
\texttt{topologization\_class\_m\_original\_complex\_platonic.tex}),
but the strict chain-level remainder-zero identity on the original
complex requires a direct Virasoro-level two-loop diagram
computation not yet carried out. This is F3 in the open-frontier
list of the Platonic reconstitution.
\end{heal}

\begin{remark}[Retract data and BRST cohomology]
\label{rem:retract-data}
The strong deformation retract
$(\pi, \iota, h)\colon \Zder(\cA) \rightleftarrows
H^{\bullet}_{Q_{\mathrm{tot}}}(\Zder(\cA))$ is not a mere bijection;
it is the data of a Cauchy--Szeg\H{o} projector for the 3d bulk BV
complex, constructed via the canonical heat-kernel regularisation on
$C\times\R_{\geq 0}$ (Costello--Gwilliam \cite{CG} Vol.~II, Appendix~A).
Transporting the $\Etop{3}$-structure from cohomology back to the
original complex via homotopy transfer
(Kontsevich--Soibelman, Lada--Markl) produces a chain-level
$\Etop{3}_{\infty}$-algebra on the zero-differential model, unique up
to $\infty$-quasi-isomorphism. This is the tier-II statement; it is
unconditional for class~L, class-G, and class-M on BRST cohomology.
\end{remark}

%=======================================================================
\section{Cycle 3: iterated Sugawara, W_N, W_infinity horizon}
\label{sec:cycle-3}
%=======================================================================

\begin{attack}[Does iterated Sugawara genuinely give $\Etop{k+2}$?]
At spin $3$, the composite field $T^{(3)}(z)$ does not close under
OPE with $T^{(2)}(z)$ in a free-field theory without $\cW_3$
corrections: the $T^{(2)} T^{(3)}$ OPE has a double pole with a
cubic primary that is not in the span of $T^{(2)}$ alone. Write out
the $\cW_3$ OPE and check that iterated Sugawara really gives
$\Etop{4}$. Chase the tower to $\Etop{\infty}$. Does the
$\cW_\infty$ limit converge operadically or is it heuristic?
\end{attack}

\begin{heal}[the Casimir tower and $\Etop{k+2}$ ladder]
\textbf{$\cW_3$ OPE algebra (Zamolodchikov 1985).} Let
$W(z) = T^{(3)}(z)$ be the spin-$3$ Casimir current. The OPE structure
is
\begin{align}\label{eq:w3-ope}
 T(z)\, W(w) &\;\sim\; \frac{3\,W(w)}{(z-w)^2}
 + \frac{\partial W(w)}{z-w}, \\
 W(z)\, W(w) &\;\sim\; \frac{c/3}{(z-w)^6}
 + \frac{2T(w)}{(z-w)^4}
 + \frac{\partial T(w)}{(z-w)^3}
 + \frac{1}{(z-w)^2}\Bigl(2\beta_3\,\Lambda(w)
 + \tfrac{3}{10}\partial^2 T(w)\Bigr) \notag \\
 & \qquad + \frac{1}{z-w}\Bigl(\beta_3\,\partial\Lambda(w)
 + \tfrac{1}{15}\partial^3 T(w)\Bigr) + \cdots,
\end{align}
with $\Lambda = \nord{TT} - (3/10)\partial^2 T$ the Zamolodchikov
composite at Virasoro norm $\langle \Lambda | \Lambda \rangle
= c(5c+22)/10$ and
$\beta_3 = 16/(22 + 5c)$. The algebra does \emph{not} close on the
Virasoro sector alone; the $W\cdot W$ OPE forces the Zamolodchikov
$\Lambda$, and closure of the iterated OPE needs the full $\cW_3$
data $(c, T, W)$ subject to the Jacobi-like identity
$[W_{(1)}, W_{(2)}] = \cdots$ determining the $\cW_3$ quantum
associativity constraint at central charge $c \neq -22/5$.

\textbf{Casimir tower axioms.}
The depth-$N$ Casimir tower is a collection of fields
$\{T^{(n)}(z)\}_{2\leq n \leq N+1}$ in $\cA$ such that:
\begin{itemize}
\item \emph{(C1)~Quasiprimarity.} Each $T^{(n)}$ is a quasiprimary
field of conformal weight $n$ with respect to the Virasoro
subalgebra $T^{(2)}$.
\item \emph{(C2)~Casimir form.} $T^{(n)} = d^{a_1\cdots a_n}\,
\nord{J_{a_1}\cdots J_{a_n}}$ for a symmetric invariant tensor
$d$ of rank $n$ on $\fg$, with $d^{ab} = g^{ab}/(k+\hv)$ at
$n = 2$ (Sugawara).
\item \emph{(C3)~BRST-primitive antighosts.} There exist
$G^{(n)}(z) \in \cA^{\mathrm{BV}}_{3d}$ of spin $n-1$ with
\begin{equation}\label{eq:Gn-BRST}
 T^{(n)}(z) \;=\; [Q_{\mathrm{tot}},\, G^{(n)}(z)]
 \qquad\text{on}\quad H^{\bullet}_{Q_{\mathrm{tot}}}.
\end{equation}
\item \emph{(C4)~Antighost commutativity.}
$[G^{(n)}(z), G^{(m)}(w)] = Q_{\mathrm{tot}}$-exact for
$2 \leq n,m \leq N+1$.
\end{itemize}

\textbf{Ladder theorem.} Let $\cA$ carry a depth-$N$ Casimir tower.
On BRST cohomology the $n$-th mode $T^{(n)}_{(0)}$ acts as
$\partial_z^{n-1}/(n-1)!$ on primary fields (Borcherds
\cite{Borcherds-vertex}, Kac \cite{Kac-vertex}). Axioms (C3)--(C4)
promote the formal flow $\{(n-1)!\,\partial_z^{n-1}\}_{n=2}^{N+1}$ on
the holomorphic formal disc to $Q_{\mathrm{tot}}$-exact. Hence on
$H^{\bullet}_{Q_{\mathrm{tot}}}(\Zder(\cA))$ the holomorphic
direction becomes $N$-fold topologised (by a Lurie-HA 5.5.3 argument
promoting locally-constant-up-to-order-$N$ to locally-constant).
Dunn additivity then gives
\[
 H^{\bullet}_{Q_{\mathrm{tot}}}(\Zder(\cA))
 \;\in\;\mathsf{Alg}_{\,\Etop{N+1}\otimes\Etop{1}}
 \;=\;\mathsf{Alg}_{\,\Etop{N+2}}.
\]
Scope: (C1)--(C3) unconditional for class~L ($V_k(\fg)$,
$k\neq -\hv$, full principal $\cW$-tower); class~M depth $N = \infty$
requires (C4) pairwise antighost commutativity, which is the
Linshaw--Zamolodchikov--Lukyanov $\cW_\infty[\mu]$ Poisson algebra on
BRST cohomology \cite{Linshaw-Winfty}; established at depth $N \leq 3$
(Fateev--Lukyanov), conditional at depth $N \geq 4$ on the
$\cW_\infty$ antighost-commutativity conjecture.

\textbf{$\cW_\infty$ horizon: convergence criterion.}
Take $N \to \infty$. The tower $\{T^{(n)}\}_{n\geq 2}$ generates
$\cW_\infty[\mu]$; the OPE algebra is determined by a single
parameter $\mu$ (the central charge is recoverable from $\mu$ and
the Gaberdiel--Gopakumar normalisation). Three independent
convergence criteria for $\cW_N \to \cW_\infty \Rightarrow
\Etop{\infty}$:
\begin{itemize}
\item \emph{(i) Axiomatic.} $\cW_\infty[\mu]$ admits a presentation
as a pro-limit of truncations $\cW_N[\mu]$ (Linshaw 2019), and
the antighost-commutativity conjecture at all $N$ promotes this
to a genuine $\Etop{\infty}$ endpoint.
\item \emph{(ii) Physical.} In the celestial holography avatar
\cite{Strominger-celestial}, the symmetry algebra of gravitational
soft theorems in $4d$ asymptotically flat gravity is
$\cW_{1+\infty}$, whose boundary avatar is the $\Etop{\infty}$
factorisation algebra. Strominger--Pasterski--Pate derive the
$w_{1+\infty}$ algebra from Einstein's equation at null infinity;
the topologisation tower gives the bulk $\Etop{\infty}$
counterpart.
\item \emph{(iii) Koszul.} The bar complex $B(\cW_\infty)$ is
$E_\infty$-coalgebraic by Dunn additivity plus the tower, and its
derived centre $\Zder(\cW_\infty)$ is $\Etop{\infty}$ on BRST
cohomology. Independent verification: the Hochschild cohomology
$HH^{\bullet}(\cW_\infty, \cW_\infty)$ vanishes in degree
$\geq 2$ on the non-logarithmic $C_2$-cofinite locus (Theorem~H),
consistent with $\Etop{\infty}$-algebra structure.
\end{itemize}
All three are independent and mutually consistent. $\cW_\infty \to
\Etop{\infty}$ is not heuristic; it is the operadic horizon of the
tower, rigorous in the $(\infty,1)$-categorical lane, conditional
on the $\cW_\infty$ antighost-commutativity conjecture at the
chain-level lane.
\end{heal}

%=======================================================================
\section{Cycle 4: the seven faces of the heptagon}
\label{sec:cycle-4}
%=======================================================================

\begin{attack}[Are seven faces seven equivalences or seven
conjectures?]
The heptagon asserts seven equivalent presentations of $\SCchtop$.
Five are classical (operadic, Koszul dual, factorisation,
BV/BRST, convolution); two are new (Drinfeld-centre,
derived-AG). List all seven precisely. Which are chain-level, which
are cohomological? Are all seven actually proved or are some
conjectural?
\end{attack}

\begin{heal}[the seven faces, explicitly]
\hfill

\begin{face}[Operadic; chain-level]\label{face:operadic}
$\SCchtop$ as a two-coloured topological operad with spaces given
by~\eqref{eq:sc-spaces}, composition via insertion of configurations,
with the directional restriction (no open-to-closed). Presentation by
codimension-one boundary strata of
$\FM_k(\mathbb{H})\times_{\Conf_m(\R)}\Conf_m(\R)$, Arnold-for-closed
plus Stasheff-for-open plus mixed codimension-two relations. Proved
Voronov 1998.
\end{face}

\begin{face}[Koszul dual; chain-level]\label{face:koszul}
$(\SCchtop)^{!} = (\Lie, \Ass, \mathrm{shuffle\text{-}mixed})$ via
Ginzburg--Kapranov bicoloured bar-cobar: closed sector $\Lie$ (from
$\Ehol{2}^{!}$), open sector $\Ass^! = \Ass$, mixed sector shuffle
symmetric, empty closed-to-open. Explicit cofibrant replacement
$W(\SCchtop) = \Omega(B(\SCchtop))$ after fixing
$\Phi\in\mathrm{GRT}_1(\Q)$.
\end{face}

\begin{face}[Factorisation; chain-level]\label{face:factorisation}
$\SCchtop$-algebras as factorisation data on
$\Ran(C)\times\Ran(\R)$: closed a factorisable $D$-module on
$\Ran(C)$ (BD chiral algebra), open a locally constant factorisation
algebra on $\Ran(\R)$ ($E_1$ by Ayala--Francis \cite{AF15}), mixed a
factorisation module supported on $C\times\{0\}$. Proved BD
\cite{BD04}, CG \cite{CG} Vol.~I Ch.~5.
\end{face}

\begin{face}[BV/BRST Swiss-cheese; chain-level]\label{face:bv}
Bulk-boundary pair in 3d holomorphic-topological BV-BRST gauge theory
on $C\times\R_{\geq 0}$. Bulk an $\Ehol{2}\otimes\Etop{1}$
factorisation algebra; boundary at $C\times\{0\}$ a chiral algebra
$\cA$; mixed the bulk-to-boundary restriction. Proved
Costello--Gaiotto \cite{CostelloGaiotto} for affine Chern-Simons
boundaries; general case in the topologisation chapter.
\end{face}

\begin{face}[Convolution / Higher Deligne;
chain-level]\label{face:convolution}
Convolution dGLA of $\Zder(\cA)$ acting on $\cA$ via braces,
$\SCchtop$-structure in brace $\infty$-operad form
\cite{Kontsevich-operads-motives}. Open colour: inner Hochschild
brace $\cA\otimes_{\cA^e}\ChirHoch(\cA,\cA)\to\cA$; closed colour:
outer Hochschild brace encoding $\Ehol{2}$. Chiral higher Deligne
theorem (Voronov-Tamarkin-Kontsevich adapted to chiral).
\end{face}

\begin{face}[Drinfeld-centre;
cohomological]\label{face:drinfeld}
Pairs $(\cA, \Drinf(\Rep_{\mathrm{fact}}(\cA)))$: closed acts on open
via $\Drinf(\Rep_{\mathrm{fact}}(\cA))\to
\mathrm{End}(\Rep_{\mathrm{fact}}(\cA))$. The Drinfeld-centre
construction is canonically a $2$-category; the $\SCchtop$-datum is
recovered at objects after Hochschild-compatible forgetful. Proved
cohomologically via Lurie HA.5.3 + factorisation centre
(Barwick--Schommer-Pries--Kan). Chain-level refinement open (F4).
\end{face}

\begin{face}[Derived-AG; cohomological]\label{face:derived-ag}
Global sections of a derived Artin stack $\mathfrak{M}_{\SCchtop}$
classifying $E_1$-chiral algebras with $\SCchtop$-datum plus its
cotangent complex. Closed: formal neighbourhood of
$\cA\in\mathfrak{M}_{E_1\text{-}\mathrm{chir}}$ in the
Hochschild-deformation direction; open: tangent $\cA$ at that point;
mixed: universal bulk-boundary infinitesimal bracket. Proved via
Lurie-Toen-Vezzosi derived deformation ($\mathrm{HA}$.5.4) +
Francis-Gaitsgory \cite{FG11}. Chain-level refinement:
$\mathfrak{M}_{\mathrm{Kosz}}$ proved on non-logarithmic $C_2$-cofinite
locus.
\end{face}

\textbf{The 5/7 chain-level vs 2/7 cohomological split.}
Faces~\ref{face:operadic}, \ref{face:koszul},
\ref{face:factorisation}, \ref{face:bv}, \ref{face:convolution} are
chain-level; \ref{face:drinfeld}, \ref{face:derived-ag} are
cohomological. The seven edges: (\ref{face:operadic}-\ref{face:koszul})
bicoloured bar-cobar; (\ref{face:koszul}-\ref{face:factorisation})
$(\Lie,\Ass)$ acting on $(D\text{-}\mathrm{mod},\mathrm{Mod})$;
(\ref{face:factorisation}-\ref{face:bv}) BD-CG equivalence (CG Thm.
5.4.9); (\ref{face:bv}-\ref{face:convolution}) Maurer-Cartan
$\leftrightarrow$ brace deformation; (\ref{face:convolution}-\ref{face:drinfeld})
chiral higher Deligne $\Rightarrow$ brace = centre;
(\ref{face:drinfeld}-\ref{face:derived-ag}) tangent to
$\mathfrak{M}_{\SCchtop}$ = Hochschild of $\cA$;
(\ref{face:derived-ag}-\ref{face:operadic}) closed-loop derived stack
$\to$ operadic. Each edge is an $\infty$-quasi-isomorphism; the
heptagon commutes. Any one face determines the $\infty$-operadic
structure; absence of a conformal vector freezes the heptagon at
$\SCchtop$, presence triggers topologisation to $\Etop{3}$, making
topologisation the exceptional collapse of the generic bicoloured
object.
\end{heal}

%=======================================================================
\section{Cycle 5: 3d quantum gravity as $\Phihol(\Vir_c)$}
\label{sec:cycle-5}
%=======================================================================

\begin{attack}[Is Virasoro Koszul? What is $\Phihol(\Vir_c)$
concretely?]
3d quantum gravity is advertised as the ``first non-trivial
instance'' of the universal holography functor, with
$\Phihol(\Vir_c)$ being a 3d HT-QFT on $C\times\R$. Virasoro is
class-M (depth $\infty$); is it Koszul in the sense required by
$\Phihol$? Write down the 3d HT-QFT concretely: its boundary
observables, its bulk observables, its partition function on
$C\times S^{1}$.
\end{attack}

\begin{heal}[$\Phihol(\Vir_c)$ and the $C\times S^1$ partition
function]
\textbf{Is Virasoro Koszul?} At non-critical $c\neq 0$ off the
$(2, 2k+1)$-minimal-model lattice, $\Vir_c$ is chirally Koszul:
$B(\Vir_c) = T^c(s^{-1}\overline{\Vir_c})$ has cohomology
concentrated in degree $0$ on the weight-completed category,
Positselski complementary duality $\Omega(B(\Vir_c)) \simeq \Vir_c$
in $D^{\mathrm{co}}_{ch}(C)$. Scope: $\mathsf{ProvedHere}$ on the
non-logarithmic $C_2$-cofinite standard landscape; class-M depth
$\infty$ only affects strictness of topologisation at chain level on
original complex (F3 frontier), not the existence of $\Phihol$.

\textbf{$\Phihol(\Vir_c)$ explicitly.} By (U1)--(U3):
$\Obs^{\partial}(\Phihol(\Vir_c)) = \Vir_c$ (boundary);
$\Obs^{\mathrm{bulk}}(\Phihol(\Vir_c)) = \Zder(\Vir_c) =
\ChirHoch(\Vir_c, \Vir_c)$ (bulk); action by $\SCchtop$-brace.
Theorem~H (Hochschild concentration) confines bulk cohomology to
degrees $0, 1, 2$; Sugawara topologisation gives the
$\Etop{3}$-structure on BRST cohomology.

\textbf{Physical identification: 3d Einstein gravity on AdS$_3$.}
Brown-Henneaux \cite{BrownHenneaux}: pure 3d gravity on AdS$_3$ has
boundary Virasoro$\times$Virasoro at $c = 3\ell/(2G_N)$, $\ell$ AdS
radius, $G_N$ Newton constant. Chern-Simons formulation (Witten
\cite{Witten3dGravity}): 3d gravity is $\mathrm{SL}_2(\R)\times
\mathrm{SL}_2(\R)$ Chern-Simons at $k = \ell/(4G_N)$; Sugawara
on $\mathrm{SL}_2$-Chern-Simons at level $k$ produces
Virasoro at $c(k) = 3k/(k+2) = 3\ell/(2G_N)$ for $k\gg 1$. Hence:
\[
 \Phihol(\Vir_c) = \text{3d HT sector of pure Einstein gravity on
 AdS}_3 \text{ at } c = 3\ell/(2G_N).
\]
Identification, not analogy: bulk $\Zder(\Vir_c)$ is local-observable
algebra of BRST-reduced 3d Chern-Simons; boundary chiral algebra is
the Liouville sector \cite{CHvD}; the $\Etop{3}$-structure on
cohomology is the topologisation of the bulk 3-form current algebra
via the inner Virasoro vector.

\textbf{Partition function on $S^2\times S^1$.}
Take $C = S^2$, compactify $\R\to S^1$ (Euclidean time,
inverse-temperature $\beta$). Modular parameter
$\tau = i\beta/(2\pi)$, $\mathrm{Im}(\tau)>0$. Partition function:
\begin{equation}\label{eq:vir-partfn}
 Z_{\Phihol(\Vir_c)}(S^2\times S^1_\tau)
 = \mathrm{Tr}_{\Vir_c}\, q^{L_0 - c/24},
 \qquad q = e^{2\pi i\tau}.
\end{equation}
For the vacuum Verma module of $\Vir_c$ at $c>1$ generic:
\begin{equation}\label{eq:vir-char}
 \chi_{\Vir_c, \mathrm{vac}}(q)
 = q^{-c/24}\prod_{n=2}^{\infty}\frac{1}{1-q^n}
 = \frac{q^{(1-c)/24}(1-q)}{\eta(\tau)}.
\end{equation}
At $c = 3\ell/(2G_N)$, this is the perturbative Brown-Henneaux
partition function of Einstein gravity on thermal AdS$_3$,
tree-level + one-loop sum (Maloney-Witten \cite{MaloneyWitten}).

\textbf{Modular-invariant partition function.} The perturbative sum
in~\eqref{eq:vir-char} is not modular-invariant; the non-perturbative
answer is obtained by summing over the $\mathrm{SL}_2(\Z)$ orbit
(Farey-tail sum, Manschot--Moore \cite{ManschotMoore}), producing
the full 3d-gravity partition function:
\begin{equation}\label{eq:vir-mod-inv}
 Z^{\mathrm{mod}}_{\Phihol(\Vir_c)}(\tau)
 \;=\; \sum_{\gamma \in
 \mathrm{SL}_2(\Z)/\mathrm{SL}_2(\Z)_{\infty}}
 \chi_{\Vir_c, \mathrm{vac}}\bigl(\gamma\cdot\tau\bigr).
\end{equation}
This is the bulk trace of the Brown--Henneaux partition function:
sum over saddle-points labelled by the Euclidean $\mathrm{SL}_2(\Z)$
black-hole geometries. At large $c$, the genus-$1$ saddle dominates
and recovers the Hardy--Ramanujan asymptotic
$\log Z \sim 2\pi\sqrt{cL_0/6}$, i.e.\ the Bekenstein--Hawking
entropy of a BTZ black hole of mass $L_0$.

\textbf{Koszul content and bar trace.}
$Z^{\mathrm{bar}}_{\Phihol(\Vir_c)}(\tau) := \mathrm{Tr}_{B(\Vir_c)}
q^{L_0 - c/24}$ equals~\eqref{eq:vir-char} by Theorem~H concentration:
bar cohomology agrees with the chiral algebra in degree $0$; degrees
$1, 2$ assemble into the $(1-q)$ prefactor. Holographic partition
function is the ordered bar trace of the boundary.

\textbf{``First non-trivial instance.''} Class-G (Heisenberg,
lattice) is abelian, trivial Sugawara; class-L (affine KM) reduces to
class-M via Sugawara, $\Phihol(V_k(\fg)) = \Phihol(\Vir_{c(k)})$
composed with Sugawara-Chern-Simons boundary; Virasoro is the first
class-M instance where $\Phihol$ captures non-trivial 3d gravity
physics with non-trivial partition function on $C\times S^1$. The
$\cW_N$ tower gives higher-spin gravitational extensions;
$\cW_\infty$ is the gravitational horizon.
\end{heal}

%=======================================================================
\section{Three-tier summary and open frontier}
\label{sec:summary}
%=======================================================================

The G3 identity is not a single theorem; it is a three-tier
architecture with one $\infty$-operadic statement and two chain-level
refinements.

\begin{center}
\begin{tabular}{|l|l|l|l|}
\hline
\textbf{Tier} & \textbf{Content} & \textbf{Scope} & \textbf{Status} \\
\hline
I & $\Phihol$ on BRST cohomology gives $\Etop{3}$ &
All class-G/L/M &
$\mathsf{ProvedHere}$ \\
II & Transferred $\Etop{3}_{\infty}$ on cohomology model &
All class-G/L/M &
$\mathsf{ProvedHere}$ \\
III & Strict chain-level $\Etop{3}$ on original complex &
Class-G/L proved; class-M open &
F3 frontier \\
\hline
\end{tabular}
\end{center}

The heptagon has 5/7 chain-level faces and 2/7 cohomological; their
mutual edges commute at both lanes where defined (chain-level edges
at the 5-face subheptagon, cohomological edges at the full heptagon).

The $\cW_N$ ladder theorem gives $\Etop{N+1}$ for depth-$N$ Casimir
towers with antighost-commutativity (C4); unconditional at $N \leq 3$,
conditional on (C4) at $N \geq 4$, and the $\cW_\infty$ horizon is
the operadic Platonic endpoint.

\textbf{Partition function on $C\times S^1$ for $\Phihol(\Vir_c)$:}
\begin{equation}\label{eq:final-partition}
 Z_{\Phihol(\Vir_c)}(S^2\times S^1_\tau) =
 \frac{q^{(1-c)/24}(1-q)}{\eta(\tau)},
 \qquad q = e^{2\pi i\tau},
\end{equation}
the Virasoro vacuum character, identified (at $c = 3\ell/(2G_N)$)
with the perturbative Brown--Henneaux partition function of 3d
Einstein gravity on thermal AdS$_3$, and its modular completion
via the Farey-tail sum over the $\mathrm{SL}_2(\Z)$ orbit.

\textbf{Cross-volume harmonies.} The G3 identity links Vol~I and
Vol~II: Vol~I supplies $\Vir_c$ as the atomic class-M example, its
$\kappa = c/2$ invariant, Koszulness at generic $c$, and the
shadow-tower depth $r_{\max} = \infty$; Vol~II carries $\Phihol$
itself and the bulk $\Zder$; at $\c = 5$ the Vol~III Mukai-enhanced
K3 Heisenberg witness refines to the Bruinier Heegner Chern-class
reciprocity output on $\cK^\kappa = 8$ (Theorem~C B-row ceiling).

\textbf{Open frontier.} (F3) Strict chain-level tier-III class-M
original-complex; (F4) pairwise antighost commutativity (C4) for
$N \geq 4$ in the Casimir tower; the $\cW_\infty$ horizon is
rigorous at tier I/II and conditional on (C4) at tier III.

%=======================================================================
\section{Synthesis: the G3 equation of state}
\label{sec:synthesis}
%=======================================================================

Assembling the five cycles:
\begin{equation}\label{eq:G3-full}
 \boxed{
 \Phihol\colon
 \ChirAlg^{\omega,\mathrm{BL}}_{C}
 \xrightarrow{\ \SCchtop\text{-braces}\ }
 \HTQFT_{C\times\R},
 \quad
 \SCchtop + \omega_{\mathrm{inner}, k\neq -\hv}
 \xrightarrow{\ \mathrm{Sug}\ } \Etop{3}
 \xrightarrow{\ \cW_N\ } \Etop{N+1}
 \xrightarrow{\ \cW_\infty\ } \Etop{\infty}.}
\end{equation}
$\SCchtop\to\Etop{3}$ is the topologisation theorem (Cycle~2);
$\Etop{3}\to\Etop{N+1}$ is the iterated-Sugawara ladder (Cycle~3);
$\Etop{\infty}$ is the $\cW_\infty$ horizon. The heptagon
$\SCchtop$ is the generic object with seven presentations
(Cycle~4); topologisation is the exceptional collapse. At
$\Phihol(\Vir_c)$ the functor realises 3d Einstein gravity on
AdS$_3$ with partition function~\eqref{eq:final-partition}
(Cycle~5).

The five-archetype behaviour of $\Phihol$: class-G
(Heisenberg/lattice) has trivially topologised bulk by ad-hoc
freeness; class-L gives affine Chern-Simons 3d gauge theory;
class-M is the non-trivial 3d HT-QFT locus with Virasoro the
first genuine instance; class-B is the Mukai-enhanced K3 Heisenberg
witness (Theorem~C B-row ceiling $\cK^\kappa = 8$ via Bruinier
Heegner Chern-class reciprocity); class-$\cW$-tower is the
$\Etop{\infty}$ celestial-gravitational horizon.

The Vol~I-Vol~II bridge at functor level:
\[
\begin{array}{ccc}
 \ChirAlg^{\omega,\mathrm{BL}}_{C} &
 \xrightarrow{\ \Phihol\ } &
 \HTQFT_{C\times\R} \\[0.5em]
 \downarrow\mathsf{av} & &
 \downarrow\pi_0 \\[0.5em]
 \mathsf{mod}\text{-}\ChirAlg_{C} &
 \xrightarrow{\ \Phihol^{\mathrm{mod}}\ } &
 \HTQFT_{C\times\R}^{\mathrm{mod}}
\end{array}
\]
$\mathsf{av}$ the ordered-to-symmetric averaging; $\pi_0$ the
zeroth-cohomology collapse of the bulk. Commutativity is the
compatibility of $E_1$-to-$E_\infty$ averaging with bulk-boundary
duality.

G3 is the statement that chiral algebra on $C$ and 3d HT-QFT on
$C\times\R$ are two presentations of one categorical datum, related
by the single functor~\eqref{eq:G3-full}.

\begin{thebibliography}{99}
\bibitem{BD04} A.~Beilinson, V.~Drinfeld, \emph{Chiral Algebras},
 AMS Colloquium Publ.~51, 2004.
\bibitem{CG} K.~Costello, O.~Gwilliam, \emph{Factorization Algebras
 in QFT, I/II}, Cambridge UP, 2017/2021.
\bibitem{Voronov1998} A.~Voronov, \emph{The Swiss-cheese operad},
 Contemp.~Math.~239 (1999), 365--373.
\bibitem{AF15} D.~Ayala, J.~Francis, \emph{Factorization homology of
 topological manifolds}, J.~Topol.~8 (2015), 1045--1084.
\bibitem{CostelloGaiotto} K.~Costello, D.~Gaiotto, \emph{VOAs and 3d
 \(\mathcal N=4\) gauge theories}, JHEP 05 (2019), 018.
\bibitem{FG11} J.~Francis, D.~Gaitsgory, \emph{Chiral Koszul duality},
 Selecta Math.~18 (2012), 27--87.
\bibitem{BrownHenneaux} J.~D.~Brown, M.~Henneaux, \emph{Central charges
 in asymptotic 3d gravity}, CMP 104 (1986), 207--226.
\bibitem{Witten3dGravity} E.~Witten, \emph{2+1 gravity as soluble
 system}, Nucl.~Phys.~B 311 (1988), 46--78.
\bibitem{CHvD} M.~Coussaert, M.~Henneaux, P.~van~Driel,
 \emph{Asymptotic 3d gravity with cosmological constant},
 CQG 12 (1995), 2961--2966.
\bibitem{MaloneyWitten} A.~Maloney, E.~Witten, \emph{3d gravity
 partition functions}, JHEP 02 (2010), 029.
\bibitem{ManschotMoore} J.~Manschot, G.~W.~Moore, \emph{A modern
 Farey tail}, CNTP 4 (2010), 103--159.
\bibitem{Strominger-celestial} A.~Strominger, \emph{Lectures on IR
 structure of gravity and gauge theory}, PUP, 2018.
\bibitem{Linshaw-Winfty} A.~Linshaw, \emph{Universal even-spin
 \(\mathcal W_\infty\)}, Adv.~Math.~366 (2020), 107079.
\bibitem{Borcherds-vertex} R.~E.~Borcherds, \emph{Vertex algebras,
 Kac-Moody, monster}, PNAS 83 (1986), 3068--3071.
\bibitem{Kac-vertex} V.~Kac, \emph{Vertex Algebras for Beginners},
 AMS ULS 10, 2nd ed.~1998.
\bibitem{Kontsevich-operads-motives} M.~Kontsevich, \emph{Operads and
 motives in deformation quantization}, LMP 48 (1999), 35--72.
\bibitem{Fateev-Lukyanov} V.~A.~Fateev, S.~L.~Lukyanov,
 \emph{\(Z_n\)-symmetric 2d CFT}, IJMPA 3 (1988), 507--520.
\end{thebibliography}

\end{document}
